(1) 用复振幅表示入射声压 \(p_i=P e^{i(kx-\omega t)}\)，对应粒子速度 \(u_i=p_i/(\rho c)\)。左行反射波满足 \(u_r=-p_r/(\rho c)\)，刚壁条件 \(u_i(0,t)+u_r(0,t)=0\)，所以压力反射系数 \(r_p=p_r(0)/p_i(0)=+1\)，而速度反射系数为 \(-1\)。
(2) 取实场，入射 \(p_i=P\cos(kx-\omega t)\)，反射 \(p_r=P\cos(kx+\omega t)\)。由和差化积得总压力 \(p=2P\cos(kx)\cos\omega t\)，总流速 \(u=2P\sin(kx)\sin\omega t/(\rho c)\)。在 \(x=0\) 速度为节点，压力为振幅 \(2P\) 的腹点，位置关系与刚壁约束一致。
(3) 入射波平均能流强度 \(I=P^2/(2\rho c)\)。声波动量通量大小为 \(I/c\)，刚壁正常完全反射使波动量法向分量反向，因此壁面所受平均辐射压力为 \(P_{\rm rad}=2I/c=P^2/(\rho c^2)\)。这属于二阶时间平均力，不能仅用一阶压强振幅代替。
(4) 在压力释放自由表面，边界为近似 \(p(0,t)=0\)，压力反射系数为 \(-1\)，与刚壁的 \(+1\) 恰相反。以上结论假定小振幅无耗散平面波、壁面尺寸远大于波长且壁固定；有限壁、非线性大声压和边界层可修正结果。
